\documentclass[12pt]{article}
\usepackage[T1]{fontenc}
\usepackage[utf8]{inputenc}
\usepackage{lmodern}
\usepackage{textcomp}
\usepackage{enumitem}
\usepackage{geometry}
\geometry{margin=1in}
\begin{document}
\begin{center}
Answer Key - History - Graduate Level Assessment 22 \
\small (Answers are ordered exactly as in the question paper.)
\end{center}

\section*{Multiple Choice}
\begin{enumerate}[label=\arabic*.]
\item \textbf{Answer:} B
\\ \textit{Options:} A) 1.8 million; 1.2 million; B) 2.5 million; 1.8 million; C) 1.4 million; 800,000; D) 4.5 million; 3.2 million; E) 2 million; 1.5 million
\\ \textit{Note:} The correct answer is based on archaeological and fossil evidence that indicates Homo habilis first appeared approximately 2.5 million years ago, with Homo erectus emerging around 1.8 million years ago, marking significant milestones in human evolution.
\item \textbf{Answer:} A
\\ \textit{Options:} A) barley, lentils, and peas.; B) soybeans, emmer, and rice.; C) millet, barley, and maize.; D) wheat, rice, and peas.; E) rice, soybeans, and maize.
\\ \textit{Note:} The excavation of Yang-shao sites reveals evidence of early agricultural practices in China, with barley, lentils, and peas being identified as key domesticated crops that contributed to the region's subsistence strategies during the Neolithic period.
\item \textbf{Answer:} A
\\ \textit{Options:} A) office; landfill; B) dining area; garbage disposal; C) bathroom; incinerator; D) kitchen; garden; E) playground; waste treatment plant
\\ \textit{Note:} The correct answer is based on the understanding that an office is a location where primary refuse, such as documents and materials generated from daily operations, is produced, while a landfill is specifically designed to contain secondary refuse, which consists of waste that has been discarded after its initial use.
\item \textbf{Answer:} A
\\ \textit{Options:} A) carbon-14 and uranium series dating; B) tephrochronology and electron spin resonance dating; C) dendrochronology and thermoluminescence dating; D) isotope-ratio mass spectrometry and fission track dating; E) lead-lead dating and uranium-thorium dating
\\ \textit{Note:} Carbon-14 and uranium series dating are correct because both methods involve measuring the decay of isotopes that have accumulated energy from their radioactive parent materials, providing a means to date archaeological artifacts and sites based on their thermal history.
\item \textbf{Answer:} B
\\ \textit{Options:} A) formal government; B) monumental earthworks; C) urban centers; D) a permanent military
\\ \textit{Note:} Monumental earthworks are a defining characteristic of the Hopewell culture, showcasing their complex social organization and ceremonial practices, even though they did not develop into a centralized state.
\end{enumerate}

\section*{True / False}
\begin{enumerate}[label=\arabic*.]
\setcounter{enumi}{5}
\item \textbf{Answer:} False
\\ \textit{Options:} A) True; B) False
\\ \textit{Note:} The correct answer is: overall the British production rose steadily throughout this period
\item \textbf{Answer:} False
\\ \textit{Options:} A) True; B) False
\\ \textit{Note:} The correct answer is: 33 1⁄3 rpm
\item \textbf{Answer:} False
\\ \textit{Options:} A) True; B) False
\\ \textit{Note:} The correct answer is: beer
\item \textbf{Answer:} False
\\ \textit{Options:} A) True; B) False
\\ \textit{Note:} The correct answer is: President Juárez
\item \textbf{Answer:} True
\\ \textit{Options:} A) True; B) False
\\ \textit{Note:} According to the passage: 12 th
\end{enumerate}

\section*{Long Form Response}
\begin{enumerate}[label=\arabic*.]
\setcounter{enumi}{10}
\item \textbf{Answer:} 27\%
\\ \textit{Note:} Expected answer: 27\%
\item \textbf{Answer:} 1945
\\ \textit{Note:} Expected answer: 1945
\end{enumerate}

\vfill
\noindent\textit{End of Answer Key}
\end{document}